\section{Introduction}
\label{sec:intro}

A world model lets a robot check an action before taking it.
Latent world models for visual control, such as DINO-WM~\cite{zhou2025dinowm} and V-JEPA~2-AC~\cite{assran2025vjepa2}, do so by predicting the features of every future frame and measuring how close the last one is to a goal image.
For choosing between candidate actions, this is both expensive and uninformative.
It is expensive because every candidate requires predicting hundreds of patch features per step, most of which repeat what the robot already sees.
It is uninformative because feature distance to a goal is a weak signal of task progress.
On PushT, even when computed on the \emph{true} future frames, it barely separates the good proposals of a diffusion policy from the bad ones: it improves success by \result{3} points, while a learned ranking of the same true futures improves it by \result{8.5}.

The opposite extreme skips prediction.
A direct scorer maps the observation, a candidate action, and the goal straight to a score~\cite{nakamoto2024vgps}.
It is cheap, but it entangles what the action does with how the outcome relates to the goal, and must learn the former anew for every goal.

We argue that a world model for decision making should predict neither every frame nor nothing, but a compact description of what the action will do: only as much as the decision needs, and only what the history does not already reveal.
The second part is the principle of conditional video coding~\cite{li2021dcvc}: when encoder and decoder share the past frames, bits are spent only on what is new.
We apply it to the future of an action that has not been executed yet.

We call the resulting model \emph{Conditional Trajectory Abstraction} (\method, \cref{fig:overview}).
During training, a target encoder compresses the actual future of an action chunk into 16 discrete tokens, conditioned on the history.
A reader scores the code against a goal image, and a feature decoder reconstructs the future from it.
Neither sees the action, so everything the action changes must pass through the code.
The world model learns to predict the code from the history and the action in a single forward pass, as the predictor of a joint-embedding architecture~\cite{lecun2022path,assran2023ijepa}.
At test time, a frozen policy proposes candidate chunks, the world model predicts one code per candidate, and the reader ranks them.
Because codes do not depend on the goal, one prediction serves any number of goals.

Our contributions are:
\begin{itemize}[leftmargin=*,topsep=2pt,itemsep=1pt]
\item a formulation of a world model's target as a conditional code of an action's future segment, with per-step latent world models and direct scorers as its two limiting cases;
\item \method, a joint-embedding world model that learns this code with a ranking reader and a conditional feature decoder, and predicts it in one pass;
\item an attribution ladder that separates what is lost by compressing the future from what is lost by predicting it;
\item experiments on PushT with a released diffusion policy, where \method \result{[result vs.\ policy, direct scorer, endpoint and per-step latent world models, with latency]}. \TODO{Second task in which progress depends on the path.}
\end{itemize}

\begin{figure*}[t]
\centering
\resizebox{\textwidth}{!}{%
\begin{tikzpicture}[
  font=\small,
  >={Stealth[length=2.2mm]},
  mod/.style={draw, rounded corners=2pt, thick, align=center, minimum height=10mm, inner sep=3pt},
  frozen/.style={mod, fill=gray!12},
  learned/.style={mod, fill=cvprblue!12, draw=cvprblue!80!black},
  ours/.style={mod, fill=cvprblue!28, draw=cvprblue!80!black, line width=1.2pt},
  data/.style={align=center, inner sep=1pt},
  tok/.style={draw, minimum size=2.2mm, inner sep=0pt, fill=cvprblue!45},
  chunk/.style={draw, minimum width=10mm, minimum height=2.2mm, inner sep=0pt, fill=orange!35},
  frame/.style={draw, fill=white, minimum width=5.5mm, minimum height=4.5mm, inner sep=0pt},
  lab/.style={font=\footnotesize, align=center, text=black!75},
  flow/.style={->, thick},
  aux/.style={->, thick, dashed, black!60}
]
% ================= deployment (top)
\node[data] (hist) at (0,0) {history $h_t$\\[1pt]\footnotesize images, proprio};
\node[frozen, anchor=west] (enc) at ($(hist.east)+(0.8,0.85)$) {frozen encoder $\varphi$\\\footnotesize DINOv2 ViT-S/14};
\node[frozen, anchor=west] (pi) at ($(hist.east)+(0.8,-0.85)$) {policy $\pi$\\\footnotesize frozen, released};
\node[data, right=7mm of enc] (C) {$C$\,{\footnotesize(cached)}};
\node[chunk, anchor=west] (a2) at ($(pi.east)+(0.7,0)$) {};
\node[chunk] (a1) at ($(a2)+(0,0.3)$) {};
\node[chunk] (aK) at ($(a2)+(0,-0.3)$) {};
\node[lab, below=0.5mm of aK] {$K$ chunks $\bA^{(k)}$};
\node[ours, anchor=west] (wm) at ({$(a2.east)+(1.3,0)$} |- hist.east) {world model $p_\theta(S\mid C,\bA)$\\\footnotesize JEPA predictor, one pass};
\begin{scope}[shift={($(wm.east)+(0.75,0)$)}]
  \foreach \y in {0.3,0,-0.3}{\foreach \c in {0,...,3}{\node[tok] at (\c*0.26,\y) {};}}
  \coordinate (codesW) at (-0.25,0);
  \coordinate (codesE) at (1.03,0);
  \coordinate (codesS) at (0.39,-0.45);
\end{scope}
\node[lab, below=0.5mm of codesS] {codes $\hat S^{(k)}$};
\node[learned, anchor=west] (reader) at ($(codesE)+(0.7,0)$) {reader $D_\psi(C,S,g)$\\\footnotesize never sees the action};
\node[data, anchor=south] (query) at ($(reader.north)+(1.0,1.05)$) {goal images $g\in\cG$};
\node[data, right=7mm of reader] (score) {$s_k$};
\node[mod, fill=white, right=7mm of score] (sel) {select $\arg\max_k s_k$\\\footnotesize execute, replan};

\draw[flow] (hist.north) |- (enc.west);
\draw[flow] (hist.south) |- (pi.west);
\draw[flow] (enc) -- (C);
\draw[flow] (pi) -- ($(a2.west)+(-0.1,0)$);
\draw[flow] (C.east) -- ++(0.3,0) |- ($(wm.west)+(0,0.25)$);
\draw[flow] ($(a2.east)+(0.1,0)$) -- ++(0.35,0) |- ($(wm.west)+(0,-0.25)$);
\draw[flow] (wm) -- (codesW);
\draw[flow] (codesE) -- (reader);
\draw[flow] (C.north) -- ++(0,0.3) -| ($(reader.north)+(-0.8,0)$);
\draw[flow] (query.south -| {$(reader.north)+(0.8,0)$}) -- ($(reader.north)+(0.8,0)$);
\draw[flow] (reader) -- (score);
\draw[flow] (score) -- (sel);
\coordinate (loopY) at (0,-2.0);
\draw[flow] (sel.south) -- (sel.south |- loopY) -- node[lab, fill=white, inner sep=1pt, pos=0.7] {observe and replan every $\He$ steps} (hist.south |- loopY) -- ($(hist.south)+(0,-0.05)$);

% ================= training (bottom), code grid aligned under the world model
\coordinate (trow) at (0,-3.8);
\begin{scope}[shift={({$(wm.center)+(-0.39,0)$} |- trow)}]
  \foreach \y in {0.3,0,-0.3}{\foreach \c in {0,...,3}{\node[tok] at (\c*0.26,\y) {};}}
  \coordinate (tcW) at (-0.25,0);
  \coordinate (tcE) at (1.03,0);
  \coordinate (tcN) at (0.39,0.45);
  \coordinate (tcS) at (0.39,-0.45);
\end{scope}
\node[lab, below=0.5mm of tcS] {code $S$};
\node[learned, anchor=east] (src) at ($(tcW)+(-0.6,0)$) {target encoder $q_\phi(S\mid C,\tau)$\\\footnotesize sees the actual future};
\foreach \i in {0,...,3}{\node[frame] (f\i) at ($(src.west)+(-2.95+\i*0.65,0)$) {};}
\node[lab, anchor=north] at ($(f1.south)!0.5!(f2.south)+(0,-0.08)$) {actual future segment $\tau$};
\node[learned, anchor=west] (rtrain) at ($(tcE)+(0.6,0)$) {reader $D_\psi$, decoder $G_\psi$\\\footnotesize see only $C$ and $S$};
\node[data, right=6mm of rtrain] (loss) {rank siblings by progress $\ell_g$\\\footnotesize reconstruct $\varphi$-features of $\tau$};
\node[lab, anchor=west, align=left, text=black, right=7mm of loss] (eqs)
{$\mathcal{L}_{\mathrm{code}}=\mathcal{L}_{\mathrm{rank}}+\alpha\,\mathcal{L}_{\mathrm{feat}}$\\[2pt]
 $\mathcal{L}_{\mathrm{WM}}=\mathcal{L}_{\mathrm{NLL}}(\sg(S))+\mathcal{L}_{\mathrm{cons}}+\mathcal{L}^{\hat S}_{\mathrm{rank}}$\\[2pt]
 \footnotesize rate bounded by $M\log_2 V=128$ bits};
\draw[flow] (f3.east) -- (src.west);
\draw[flow] (src) -- (tcW);
\draw[flow] (tcE) -- (rtrain);
\draw[flow] (rtrain) -- (loss);
\draw[aux] (tcN) -- node[lab, right, fill=black!3, inner sep=1pt, pos=0.35] {target for the world model (stop-gradient)} (wm.south);
\begin{scope}[on background layer]
  \node[draw=black!30, fill=black!3, rounded corners=3pt, fit=(f0)(src)(eqs)(tcS)(tcN), inner xsep=4mm, inner ysep=5mm] (tbox) {};
\end{scope}
\node[lab, anchor=north west, font=\footnotesize\bfseries, text=black] at ($(tbox.north west)+(1.5mm,-1mm)$) {Training only};
\node[lab, anchor=south west, font=\footnotesize\bfseries, text=black] at ($(hist.north west)+(-0.2,1.5)$) {Deployment: predict once, query many};
\end{tikzpicture}}
\caption{\textbf{Conditional Trajectory Abstraction.}
\emph{Top:} at each decision, the history is encoded once into a context $C$ of frozen DINOv2 features, and a frozen policy proposes $K$ action chunks.
In one batched pass, the world model predicts for each chunk a 16-token code $\hat S$ of the chunk's entire future segment, conditioned on $C$.
A reader that sees $C$ and the code, but never the action, scores each code against goal images; one prediction serves any number of goals.
\emph{Bottom:} during training, a target encoder compresses the actual future segment into the code given $C$.
The reader ranks sibling candidates by progress, and a decoder reconstructs the future's frozen features from $(C,S)$.
The world model predicts the stop-gradient code and is also trained through the frozen reader.}
\label{fig:overview}
\end{figure*}

